\section{Numerical realization: a symmetric B-spline potential}
\subsection{Construct even basis functions rather than repair fitted maps}
Scale coordinates to $x=i_d/I_*$ and $y=i_q/I_*$, and approximate
$\coenergy=\psi_*I_*\normalizedcoenergy$. On symmetric q knots, define paired basis functions
\begin{equation}
 E_n(y)=\tfrac12\bigl(B_n(y)+B_n(-y)\bigr),\qquad
 \normalizedcoenergy(x,y)=\sum_{m,n}c_{mn}B_m(x)E_n(y).\label{eq:spline-potential}
\end{equation}
Keep one representative of each reflected pair to avoid duplicated columns.
Then $E_n$ is even, $E'_n$ is odd, and $E''_n$ even. Analytic differentiation
gives
\begin{align}
 \psi_d/\psi_*&=\sum c_{mn}B'_mE_n,&
 \psi_q/\psi_*&=\sum c_{mn}B_mE'_n,\\
 L_{dd}/L_*&=\sum c_{mn}B''_mE_n,&
 L_{dq}/L_*&=L_{qd}/L_*=\sum c_{mn}B'_mE'_n,\\
 L_{qq}/L_*&=\sum c_{mn}B_mE''_n,&L_*&=\psi_*/I_*.
\end{align}
The structural properties now hold by construction, including between samples.
Differential inductance is surface curvature; its change across the current
plane expresses saturation and cross-saturation.

For degree $p$ with simple interior knots a spline is $C^{p-1}$ there.
A cubic potential therefore has continuous Hessian but its third derivatives
can jump at knots. The illustration uses degree four to obtain continuous
third derivatives at simple interior knots. Repeated knots reduce continuity;
the intended derivative quality must guide knot multiplicities. Local support
limits the number of active coefficients per sample. The standard B-spline
definition and analytic derivatives are documented in \cite{scipy-bspline}.

\subsection{A linear gradient design and regularization}
Each sample adds two rows $[B'_mE_n]$ and $[B_mE'_n]$ to a design matrix
$\matr G$. The normalized data stack is $\vect z$. With a roughness operator
$\matr P$, solve
\begin{equation}
 \min_{\vect c}\norm{\matr G\vect c-\vect z}^2
          +\lambda\norm{\matr P\vect c}^2,
 \qquad \vect g_0\trans\vect c=0.\label{eq:spline-solve}
\end{equation}
Here $\vect g_0$ evaluates the potential at the reference. A null-space basis
of the gauge row gives an unconstrained least-squares problem in fewer
variables. QR or SVD exposes residual rank deficiency without forming a
poorly conditioned normal matrix. For the synthetic experiment $\matr P$
samples $\normalizedcoenergy_{xxx},\sqrt3\normalizedcoenergy_{xxy},\sqrt3\normalizedcoenergy_{xyy},\normalizedcoenergy_{yyy}$ on a uniform grid,
with quadrature weights. It approximates the squared norm of the third
derivative tensor and penalizes changing Hessian rather than constant inductance.

\subsection{Alternative representations and what they guarantee}
Independent unconstrained RBF surfaces can accommodate scattered data but
enforce neither parity nor integrability by default. Independent even/odd
flux bases enforce parity but leave the mixed derivatives unrelated.
A scalar symmetry-aware RBF potential fitted through its gradients enforces
both, like the spline potential, if the basis is sufficiently differentiable.
RBF kernel shape, smoothing, polynomial augmentation, conditioning, and
support all require tuning; the standard interpolation interface documents
these numerical choices \cite{scipy-rbf}. No universal optimal basis is claimed.

\begin{table}[tb]
\centering\small
\begin{tabularx}{\textwidth}{@{}lYY@{}}
\toprule
Representation & Structural guarantee & Remaining numerical questions\\
\midrule
Independent flux RBFs & None by default & Kernel/shape, smoothing, conditioning,
derivatives, boundary support\\
Independent parity fits & d-even and q-odd & Reciprocity, derivative quality,
sample coverage, regularization\\
Common symmetric potential & Parity and reciprocity for a smooth basis &
Gradient-design rank, gauge, bias, curvature, convexity, tuning\\
\bottomrule
\end{tabularx}
\caption{Structural comparison. Guarantees do not establish which method has
the smallest prediction error on a particular measured data set.}
\end{table}
The small numerical comparison below uses B-splines for all three structural
classes to isolate constraints from kernel selection. RBF comparison remains
conceptual in this draft; a broad interpolator benchmark is separate work.
